\documentclass[tikz,border=5pt]{standalone}
\usepackage{ctex}
\usepackage{amsmath,amssymb}
\usetikzlibrary{arrows.meta,positioning,decorations.pathreplacing}
% Shared visual language for LFS architecture diagrams.
\RequirePackage{../../mymath}

\definecolor{LFSBlue}{HTML}{2563EB}
\definecolor{LFSBlueSoft}{HTML}{DBEAFE}
\definecolor{LFSGreen}{HTML}{15803D}
\definecolor{LFSGreenSoft}{HTML}{DCFCE7}
\definecolor{LFSOrange}{HTML}{C2410C}
\definecolor{LFSOrangeSoft}{HTML}{FFEDD5}
\definecolor{LFSRed}{HTML}{B91C1C}
\definecolor{LFSRedSoft}{HTML}{FEE2E2}
\definecolor{LFSPurple}{HTML}{7E22CE}
\definecolor{LFSPurpleSoft}{HTML}{F3E8FF}
\definecolor{LFSGray}{HTML}{475569}
\definecolor{LFSGraySoft}{HTML}{F1F5F9}

\tikzset{
  lfs picture/.style={
    font=\small,
    node distance=8mm and 12mm,
    >=Latex,
    every path/.style={line cap=round, line join=round}
  },
  block/.style={
    draw=LFSBlue,
    fill=LFSBlueSoft,
    rounded corners=2pt,
    line width=0.8pt,
    align=center,
    inner xsep=7pt,
    inner ysep=5pt,
    minimum height=10mm
  },
  compute/.style={
    block,
    draw=LFSPurple,
    fill=LFSPurpleSoft
  },
  storage/.style={
    block,
    draw=LFSGreen,
    fill=LFSGreenSoft
  },
  interface/.style={
    block,
    draw=LFSOrange,
    fill=LFSOrangeSoft
  },
  risk/.style={
    block,
    draw=LFSRed,
    fill=LFSRedSoft
  },
  neutral/.style={
    block,
    draw=LFSGray,
    fill=LFSGraySoft
  },
  decision/.style={
    diamond,
    draw=LFSOrange,
    fill=LFSOrangeSoft,
    line width=0.8pt,
    align=center,
    inner sep=2pt,
    aspect=2
  },
  group/.style={
    draw=LFSGray,
    rounded corners=3pt,
    dashed,
    line width=0.7pt,
    inner sep=6pt
  },
  arrow/.style={
    -{Latex[length=2.4mm,width=1.6mm]},
    draw=LFSGray,
    line width=0.9pt
  },
  data arrow/.style={
    arrow,
    draw=LFSBlue
  },
  control arrow/.style={
    arrow,
    draw=LFSOrange,
    dashed
  },
  residual/.style={
    arrow,
    draw=LFSGreen,
    rounded corners=4pt
  },
  label text/.style={
    font=\scriptsize,
    text=LFSGray,
    fill=white,
    inner sep=1.5pt
  },
  title/.style={
    font=\bfseries,
    text=LFSGray,
    align=center
  }
}


\begin{document}
\begin{tikzpicture}[lfs picture, >=Latex]

% Panel 1: Top-k Truncation (Fixed k)
\begin{scope}[xshift=0cm]
  \node[font=\small\bfseries, text=LFSGray] at (2.5, 4.2) {Top-$k$ 截断采样 (固定数量 $k=4$)};

  % Axes
  \draw[->, line width=0.6pt] (0, 0) -- (5.2, 0) node[right, font=\tiny] {词元降序排名};
  \draw[->, line width=0.6pt] (0, 0) -- (0, 3.8) node[above, font=\tiny] {词元概率 $p_i$};

  % Bars for Top-k
  % Selected (1 to 4)
  \fill[LFSBlueSoft, draw=LFSBlue, line width=0.6pt] (0.2, 0) rectangle ++(0.5, 3.2);
  \node[font=\tiny, above] at (0.45, 3.2) {$0.38$};
  \node[font=\tiny, text=LFSGray, below] at (0.45, -0.05) {$v_{(1)}$};

  \fill[LFSBlueSoft, draw=LFSBlue, line width=0.6pt] (0.9, 0) rectangle ++(0.5, 2.3);
  \node[font=\tiny, above] at (1.15, 2.3) {$0.27$};
  \node[font=\tiny, text=LFSGray, below] at (1.15, -0.05) {$v_{(2)}$};

  \fill[LFSBlueSoft, draw=LFSBlue, line width=0.6pt] (1.6, 0) rectangle ++(0.5, 1.4);
  \node[font=\tiny, above] at (1.85, 1.4) {$0.16$};
  \node[font=\tiny, text=LFSGray, below] at (1.85, -0.05) {$v_{(3)}$};

  \fill[LFSBlueSoft, draw=LFSBlue, line width=0.6pt] (2.3, 0) rectangle ++(0.5, 0.9);
  \node[font=\tiny, above] at (2.55, 0.9) {$0.10$};
  \node[font=\tiny, text=LFSGray, below] at (2.55, -0.05) {$v_{(4)}$};

  % Cutoff line at k=4
  \draw[dashed, draw=LFSRed, line width=1pt] (3.0, -0.2) -- (3.0, 3.6);
  \node[font=\scriptsize\bfseries, text=LFSRed, above] at (3.0, 3.6) {截断边界 $k=4$};

  % Truncated tails
  \fill[LFSGray!20, draw=LFSGray!60, line width=0.5pt] (3.2, 0) rectangle ++(0.4, 0.4);
  \fill[LFSGray!20, draw=LFSGray!60, line width=0.5pt] (3.8, 0) rectangle ++(0.4, 0.25);
  \fill[LFSGray!20, draw=LFSGray!60, line width=0.5pt] (4.4, 0) rectangle ++(0.4, 0.15);

  \node[font=\tiny, text=LFSRed, align=center] at (4.0, 1.5) {
    被截断长尾\\
    概率置零
  };

  \node[font=\scriptsize, text=LFSGray, align=center] at (2.5, -0.8) {
    无论分布陡峭或平坦，严格保留前 $k$ 项；\\
    在平坦分布下可能误杀高概率候选项。
  };
\end{scope}

% Separator
\draw[dashed, draw=LFSGray!40, line width=0.5pt] (5.8, -1.2) -- (5.8, 4.4);

% Panel 2: Top-p Nucleus Truncation (Adaptive)
\begin{scope}[xshift=6.5cm]
  \node[font=\small\bfseries, text=LFSGray] at (2.6, 4.2) {Top-$p$ 核采样 (自适应累积质量 $\rho=0.85$)};

  % Axes
  \draw[->, line width=0.6pt] (0, 0) -- (5.4, 0) node[right, font=\tiny] {词元降序排名};
  \draw[->, line width=0.6pt] (0, 0) -- (0, 3.8) node[above, font=\tiny] {累积概率 $\sum p_i$};

  % Cumulative curve
  \draw[line width=1.2pt, draw=LFSPurple]
    (0.4, 0.8) -- (1.1, 1.8) -- (1.8, 2.5) -- (2.5, 2.9) -- (3.2, 3.15) -- (4.0, 3.3) -- (4.8, 3.4);

  % Threshold line rho = 0.85 (scaled to y=2.8)
  \draw[dashed, draw=LFSGreen, line width=1pt] (0, 2.8) -- (5.0, 2.8);
  \node[font=\scriptsize\bfseries, text=LFSGreen, above right] at (0, 2.8) {累积阈值 $\rho=0.85$};

  % Cutoff marker
  \filldraw[fill=LFSRedSoft, draw=LFSRed, line width=0.8pt] (2.5, 2.9) circle (2.5pt);
  \draw[dotted, draw=LFSRed, line width=0.8pt] (2.5, 0) -- (2.5, 2.9);
  \node[font=\tiny, text=LFSRed, below] at (2.5, -0.05) {$m=3$ (核集合)};

  % Nucleus set bracket
  \draw[decorate, decoration={brace, amplitude=4pt, mirror}, line width=0.7pt, draw=LFSPurple]
    (0.2, -0.2) -- (2.5, -0.2);
  \node[font=\tiny\bfseries, text=LFSPurple, below=4pt] at (1.35, -0.2) {核集合 $\mathcal{N}_\rho$};

  \node[font=\scriptsize, text=LFSGray, align=center] at (2.6, -1.0) {
    动态根据累积概率选择最小前缀集合；\\
    高置信时动态收缩至极少数词元，\\
    低置信时自适应扩展候选范围。
  };
\end{scope}

\end{tikzpicture}
\end{document}
